Alice 点击购买后，页面转了几秒，显示连接超时。她没有看到订单号，也没有看到扣款号。她又点了一次，这回页面很快显示购买成功。打开账户明细，却有两条各为 30.00 元的扣款。

两次扣款各自对应一张订单和一条流水，余额也与流水合计相符。服务收到两次 HTTP 请求，分别完成了两笔购买。Alice 第二次点击时，只想继续完成刚才的购买；服务保存的记录中，却没有这两次点击之间的关系。

一次购买已经提交，与客户端知道它已经提交，中间隔着一次结果传递。响应在这里丢失，客户端就得在信息不全时决定下一步。重试、查询、等待和取消分别能取得哪些信息、改变哪些状态，是本章要沿着这次扣款逐步分析的问题。购买流程还需要保存一种跨请求的操作身份：让第一次发送和以后为完成同一购买而发出的请求，都能指向 Alice 授权的这一次购买。
